% =============================================================================
% Modern Back Cover - CORRETTA E OTTIMIZZATA PER 15x21cm
% =============================================================================

\pagestyle{empty}
\cleardoublepage

\begin{tikzpicture}[remember picture,overlay]
	
	% 1. SFONDO (Tutto nero/modernblack)
	\fill[modernblack] (current page.south west) rectangle (current page.north east);
	
	% 2. CITAZIONE IN ALTO (Anchor: top center)
	% Posizionata a 1.5cm dal bordo superiore per guadagnare spazio verticale
	\node[name=topquote, anchor=north, text width=0.8\textwidth, align=center] 
	at ([yshift=-1.5cm] current page.north) {%
		\color{white}
		{\sffamily\Large\textit{%
				``Il futuro è già qui,\\
				solo non è distribuito uniformemente.''
		}}\\[0.5em]
		{\sffamily\normalsize\textls[50]{WILLIAM GIBSON}}
	};
	
	% 3. DESCRIZIONE CENTRALE
	% Ancorata SOTTO la citazione (topquote.south) con 0.8cm di distacco.
	% Usiamo \small per far stare tutto nel formato 15x21.
	\node[anchor=north, text width=0.8\textwidth, align=justify] 
	at ([yshift=-0.8cm] topquote.south) {%
		\sffamily\small\color{white} % Font leggermente più piccolo e bianco
		\setlength{\parskip}{0.8em}  % Spazio tra i paragrafi
		
		Viviamo nell'era dell'informazione con cervelli progettati per la savana. 
		Questo libro esplora il paradosso fondamentale della condizione umana moderna: 
		possediamo la tecnologia più avanzata della storia, ma la gestiamo con menti 
		forgiate da milioni di anni di evoluzione per tutt'altri scopi.
		
		Dalle trappole cognitive dei social media alle decisioni irrazionali negli 
		investimenti, dal sovraccarico informativo alla dipendenza da notifiche, 
		\textit{\booktitle} svela i meccanismi nascosti che 
		governano il nostro comportamento digitale quotidiano.
		
		Un viaggio illuminante attraverso neuroscienze, psicologia evolutiva e 
		tecnologia che cambierà il modo in cui vedete voi stessi e il mondo 
		digitale che vi circonda.
	};
	
	% 4. ISBN / BARCODE (Angolo in basso a destra)
	% Il box bianco per il codice a barre
	\node[fill=white, minimum width=1.5in, minimum height=0.8in, anchor=south east] 
	at ([xshift=-1cm, yshift=1cm] current page.south east) {%
		\color{modernblack}\tiny
		\begin{tabular}{c}
			ISBN \isbn\\[0.2em]
			% Scommenta la riga sotto quando avrai l'immagine del barcode
			%\includegraphics[width=1.3in, height=0.5in]{barcode-placeholder.png}
		\end{tabular}
	};
	
	% 5. PREZZO (Angolo in basso a sinistra)
	% Allineato verticalmente con il box del barcode
	\node[white, anchor=south west] 
	at ([xshift=1cm, yshift=1.2cm] current page.south west) {%
		{\sffamily\small €\,18,00}
	};
	
	% 6. TESTO SUL DORSO/LATO (Opzionale)
	% Ruotato di 90 gradi
	\node[white, rotate=90, anchor=south] 
	at ([xshift=-0.5cm] current page.west) {%
		{\sffamily\tiny\textls[100]{SAGGISTICA \textbullet{} NEUROSCIENZE}}
	};
	
\end{tikzpicture}